\section{Dataset}\label{sec:dataset}

\subsection{Fuente oficial}
Se utilizó la base \emph{SINIESTROS DE TRANSITO FATALES 2021--2025 (PRELIMINAR)}, publicada el 2026-02-27 por el Observatorio Nacional de Seguridad Vial (ONSV) en su portal de datos abiertos:

\begin{quote}
\url{https://www.onsv.gob.pe/datosabiertos}
\end{quote}

\noindent El portal describe el recurso como: ``Detalle de siniestros de tránsito con consecuencias fatales, ocurridos a nivel nacional, 2021--2025 (preliminar)'', autor ONSV. El archivo original se mantiene en \texttt{data/raw/BBDD\_ONSV\_SINIESTROS\_FATALES\_2021-2025.xlsx} (hoja \texttt{SINIESTROS}, encabezado en la fila 5). La fecha de descarga y verificación local fue el 2026-07-09. Es un dato abierto del Estado peruano, lo que permite libre reuso con atribución. El ONSV declara la información de 2025 como preliminar.

\subsection{Por qué se cambió la fuente}
La primera iteración del proyecto usó el registro Sutran 2020--2021 de accidentes en carreteras, que permite construir el target mortal/no-mortal pero solo ofrece seis variables base (fecha, hora, departamento, vía, kilómetro y modalidad), sin ninguna condición del entorno. La evaluación crítica de esa iteración mostró un techo informacional: la red empataba con la regresión logística porque no había señal suficiente que aprender. La base ONSV aporta las variables pre-impacto que la literatura reconoce como determinantes (clima, geometría y superficie de la vía, ruralidad, red vial, coordenadas), al costo de restringir el universo a siniestros fatales. Se aceptó ese costo reformulando el problema hacia la alta letalidad, en lugar de mezclar fuentes con criterios de registro incompatibles.

\subsection{Diccionario de datos base}
La \autoref{tab:diccionario} resume las variables del registro y su rol en el proyecto.

\begin{table}[H]
\centering
\caption{Variables del registro ONSV y su rol en el proyecto.}
\label{tab:diccionario}
\begin{tabular}{lll}
\toprule
Columna & Rol & Justificación \\
\midrule
\texttt{FECHA}, \texttt{HORA} & Característica (derivar) & Estacionalidad, nocturnidad, fatiga \\
\texttt{DEPARTAMENTO} & Característica categórica & Geografía del riesgo \\
\texttt{ZONA} & Característica categórica & Ruralidad: velocidades y respuesta médica \\
\texttt{RED VIAL} & Característica categórica & Nacional/departamental/provincial/urbana \\
\texttt{TIPO DE VÍA} & Característica categórica & Carretera, avenida, calle, vía expresa \\
\texttt{COD CARRETERA} & Característica (derivar) & Familias de vías y frecuencia \\
\texttt{COORDENADAS} & Característica numérica & Señal espacial continua \\
\texttt{CONDICIÓN CLIMÁTICA} & Característica categórica & Visibilidad y adherencia \\
\texttt{CARACTERÍSTICAS DE VÍA} & Característica categórica & Recta, curva, intersección, estructura \\
\texttt{PERFIL LONGITUDINAL} & Característica categórica & Plana vs.\ inclinada \\
\texttt{SUPERFICIE DE CALZADA} & Característica categórica & Pavimentada, afirmado, trocha \\
\texttt{CLASE SINIESTRO} & Característica categórica & Mecánica del evento \\
\texttt{CANTIDAD DE FALLECIDOS} & \textbf{Solo target} & \textbf{Prohibida como entrada} \\
\texttt{CANTIDAD DE LESIONADOS} & \textbf{Excluida} & Resultado posterior (fuga) \\
\texttt{CANTIDAD VEHÍCULOS DAÑADOS} & \textbf{Excluida} & Resultado posterior (fuga) \\
\texttt{CAUSA (factor y específica)} & \textbf{Excluida} & Se determina en la investigación posterior \\
\texttt{SEÑALES vertical/horizontal} & \textbf{Excluida} & 78\,\% de faltantes concentrados por periodo \\
\bottomrule
\end{tabular}
\end{table}

\subsection{Variable objetivo}
Se construyó el target binario condicional al universo de siniestros fatales:
\[
y = 1 \iff \texttt{FALLECIDOS} \geq 2 \quad \text{(``ALTA LETALIDAD'')}, \qquad y = 0 \iff \texttt{FALLECIDOS} = 1
\]
La prohibición de \texttt{LESIONADOS}, \texttt{VEHICULOS\_DANADOS} y las columnas de causa como entradas se declara explícitamente como decisión metodológica en la \autoref{sec:preprocesamiento}: los conteos se levantan después del siniestro (fuga directa del resultado) y la causa se codifica al cierre de la investigación (el 49\,\% figura ``en proceso de investigación'').

\subsection{Auditoría inicial}
La auditoría (Bloque B) verificó y registró:
\begin{itemize}
  \item Formato: Excel (\texttt{.xlsx}), hoja \texttt{SINIESTROS}, encabezado en fila 5, 27 columnas.
  \item Número de filas de datos: 9\,106.
  \item Fechas inválidas eliminadas: 2. Duplicados por código de siniestro: 0.
  \item N final útil: 9\,104.
  \item Casos multifatales: 928 (10.19\,\%); letalidad simple: 8\,176.
  \item Rango temporal: 2021-01-01 a 2025-12-30 (2025 parcial, 678 registros).
  \item Coordenadas fuera del territorio peruano anuladas: 1 registro.
\end{itemize}

Como $3\,000 \leq N < 10\,000$ se aplicó la contingencia C2 (arquitectura A, \emph{split} 70/15/15). Como el porcentaje de clase multifatal se ubica entre 5\,\% y 25\,\%, se aplicó C5: \texttt{class\_weight} sin SMOTE. \autoref{tab:audit} resume la auditoría.

\begin{table}[H]
\centering
\caption{Resumen de la auditoría del dataset (Bloque B).}
\label{tab:audit}
\begin{tabular}{lr}
\toprule
Indicador & Valor \\
\midrule
Filas de datos originales & 9\,106 \\
Fechas inválidas eliminadas & 2 \\
Duplicados por código & 0 \\
Filas útiles finales & 9\,104 \\
Casos multifatales (2+ fallecidos) & 928 \\
Porcentaje de clase multifatal & 10.19\,\% \\
Rango temporal & 2021--2025 (2025 parcial) \\
Máximo de fallecidos en un siniestro & 33 \\
\bottomrule
\end{tabular}
\end{table}